% !TeX program = lualatex
% !TeX encoding = utf8
% !TeX spellcheck = uk_UA
% !TeX root =../PyscfBook.tex

%% =========================================================
\section{Прості атомні системи}
%% =========================================================

Прості атомні системи --- це найзручніша відправна точка для вивчення методів квантової хімії та самозгодженої процедури поля (SCF).
На прикладі атомів Гідрогену та Гелію можна продемонструвати основні ідеї методу Гартрі–Фока, вплив базисного набору, збіжність енергії та базові аспекти орбітального аналізу.

%% --------------------------------------------------------
\subsection{Атом Гідрогену}
%% --------------------------------------------------------

Атом Гідрогену є найпростішим квантовим об'єктом: він складається лише з одного електрона, що рухається в кулонівському полі ядра. Його гамільтоніан має вигляд
\[
\hat{H} = -\frac{1}{2}\nabla^2 - \frac{1}{r},
\]
де перший доданок описує кінетичну енергію електрона, а другий --- потенціал притягання до ядра.
Ця система має аналітичне розв’язання, і енергія основного стану відома точно:
\[
E_1 = -\frac{1}{2}\ \text{Ha}.
\]

Метод Гартрі–Фока для одноелектронних систем не містить апроксимацій --- відсутнє електрон–електронне відштовхування, отже HF-енергія є точною для будь-якого заданого базису. Вся похибка зумовлена лише обмеженістю базисного набору.

\inputcode[mode=inline]{H_SCF_calculation.py}

%Для мінімального базису \texttt{STO-3G} отримується енергія близька до \(-0.5\) Ha із похибкою близько \(10^{-3}\) Ha.
%Цей приклад зручний для демонстрації того, як базис впливає на точність квантово-хімічних розрахунків.
%
%%% --------------------------------------------------------
%\subsubsection{Залежність від базисного набору}
%%% --------------------------------------------------------

Щоб побачити, як збільшується точність при розширенні базису, порівняймо кілька популярних наборів --- від \texttt{STO-3G} до \texttt{cc-pV5Z}. Код знаходиться в файлі \inputcode[]{h_atom_basis_convergence.py}, а графік, який виводить код подано на (рис.~\ref{pic:h_atom_basis_convergence}).

%Для кожного базису обчислюється енергія та похибка відносно точного значення \(-0.5\) Ha.


%---------------------------------------------------------
\begin{figure}[h!]\centering
\includegraphics[width=\linewidth]{\currfiledir pictures/h_atom_basis_convergence.pdf}
\caption{Залежність енергії атома \ce{H} від вибору базису.}
\label{pic:h_atom_basis_convergence}
\end{figure}
%---------------------------------------------------------

З графіка (рис.~\ref{pic:h_atom_basis_convergence}) видно, що при переході до великих кореляційно-узгоджених базисів (наприклад, \texttt{cc-pVQZ}, \texttt{cc-pV5Z}) енергія швидко збігається до повного базисного ліміту (CBS limit).

%% --------------------------------------------------------
\subsubsection{Орбіталі та матриця густини}
%% --------------------------------------------------------

Попри те, що в атома Гідрогену є лише одна заповнена орбіталь (1s), програма \texttt{PySCF} дозволяє аналізувати всі орбіталі базису, енергетичний спектр і матрицю густини.

\inputcode[mode=inline]{uhf_h_atom.py}

\begin{minted}{text}
Енергія 1s орбіталі: -0.49980981 Ha
Теоретична енергія:  -0.50000000 Ha
Слід dm(α): 0.5
\end{minted}

Слід матриці густини для альфа-спіну дорівнює \(0.5\), оскільки \texttt{PySCF} нормує густину на повну систему (\(\alpha+\beta\)).
Для одноелектронного випадку:
\[
\mathrm{Tr}(\mathrm{dm}_\alpha) = 0.5, \quad \mathrm{Tr}(\mathrm{dm}_\beta) = 0, \quad
\mathrm{Tr}(\mathrm{dm}_\alpha + \mathrm{dm}_\beta) = 1.
\]

Отриманий результат ідеально відтворює точний розв'язок, ілюструючи коректність SCF-процедури для одноелектронних систем.

%% --------------------------------------------------------
\subsection{Атом Гелію}
%% --------------------------------------------------------

Атом гелію --- найпростіша багатoелектронна система і перший приклад, де проявляється \textbf{електрон–електронна взаємодія}.
Обидва електрони заповнюють орбіталь $1s$ з протилежними спінами, утворюючи замкнену оболонку (синглет, $S=0$).

\inputcode[mode=inline]{HeHF.py}

Після збіжності SCF-розрахунку можна отримати компоненти енергії з одно- та двоелектронних інтегралів і перевірити віральне співвідношення:
\[
\frac{\langle V \rangle}{\langle T \rangle} \approx -2.00.
\]
Відхилення від цього значення свідчить про неточність базису або проблеми збіжності.

%% --------------------------------------------------------
\subsubsection{Порівняння RHF та UHF}
%% --------------------------------------------------------

Оскільки гелій має рівну кількість електронів обох спінів (\(N_\alpha = N_\beta\)), обидва варіанти методу --- \texttt{RHF} і \texttt{UHF} --- дають однакові результати.

\inputcode[mode=inline]{rhf_vs_uhf_he.py}

\begin{commentbox}
Для відкрито-оболонкових систем метод UHF може порушувати спінову симетрію (\textit{spin contamination}).
У випадку гелію цього не відбувається: $S^2 = 0$, тому RHF $\equiv$ UHF.
\end{commentbox}

%% --------------------------------------------------------
\subsection{Збуджені стани в межах методу Хартрі–Фока}
%% --------------------------------------------------------

Метод Хартрі–Фока у своїй стандартній формі призначений для опису
\textbf{основного стану} системи — стану з найменшою можливою енергією.
Під час самозгодженої процедури (SCF) електрони заповнюють орбіталі
починаючи з найнижчих енергій, і програма автоматично знаходить конфігурацію,
що мінімізує повну енергію:
\[
E_0 = \min_{\Phi} \langle \Phi | \hat{H} | \Phi \rangle,
\]
де $\Phi$ — це одно-детермінантна хвильова функція (детермінант Слейтера).
Отже, звичайний HF-розрахунок завжди повертає \emph{основний стан}.

Для опису \textbf{збуджених станів} потрібно враховувати мультиплетність системи
та іноді примусово фіксувати електронну конфігурацію:

\textbf{Триплетні стани.} Для атома гелію збуджений триплетний стан $1s2s\,^3S$
можна отримати простим встановленням параметра \inlinecode{spin = 2}, що відповідає
двом неспареним електронам. Оскільки триплетна конфігурація вимагає, щоб електрони
мали паралельні спіни і знаходилися на різних орбіталях, SCF-процедура автоматично
знаходить конфігурацію $1s2s$ без додаткових втручань.

\textbf{Синглетні збуджені стани.} Для збудженого синглету $1s2s\,^1S$ ситуація складніша.
Параметр \inlinecode{spin = 0} означає лише, що кількість альфа- та бета-електронів
однакова, але не визначає, \emph{які саме орбіталі} вони займають. Тому звичайна
SCF-процедура знову поверне основний стан $1s^2$, оскільки він енергетично вигідніший.

Щоб отримати збуджений синглет, необхідно \textbf{зафіксувати електронну конфігурацію}
і не дозволити системі сповзти до основного стану. Для цього використовується
\textbf{метод MOM (Maximum Overlap Method)}, який на кожній ітерації SCF обирає орбіталі
з найбільшим перекриттям з орбіталями попередньої ітерації:
\[
O_{ij} = \left| \langle \phi_i^{(n)} \mid \phi_j^{(n-1)} \rangle \right|^2.
\]
Завдяки цьому MOM утримує бажану конфігурацію (наприклад, $1s2s$) і не дозволяє
хвильовій функції сповзати до основного стану ($1s^2$).

Отриманий результат називають \textbf{самозгодженим збудженим детермінантом}
(self-consistent excited determinant). Такий детермінант \emph{не є справжнім власним
станом гамільтоніана} $\hat{H}$, оскільки він не задовольняє рівняння Шредінґера точно:
\[
\hat{H}\Phi \ne E\Phi.
\]
Метод Хартрі–Фока обмежує хвильову функцію лише одним детермінантом, тоді як
справжній власний стан гамільтоніана --- це, взагалі кажучи, комбінація багатьох
детермінантів (багатоконфігураційна хвильова функція). Тому \textbf{енергії
збуджених станів, як і основного, систематично завищені} порівняно з експериментом.

Проте такі стани мають практичну цінність. Їх енергії та орбіталі дозволяють:
\begin{itemize}
    \item побачити, \textbf{де локалізований збуджений електрон} — наприклад, у гелію
    при переході $1s \to 2s$ один електрон залишається поблизу ядра, а інший
    зміщується далі, змінюючи розподіл електронної густини $\rho(\mathbf{r})$;
    \item оцінити \textbf{енергетичні інтервали між станами}
    \[
    \Delta E = E_{\text{збудж}} - E_{\text{основ}},
    \]
    які відповідають спостережуваним спектральним лініям у дослідах поглинання чи випромінювання.
\end{itemize}

Приклад реалізації MOM-розрахунку для атома гелію наведено у файлі
\inputcode{he_excited_states_mom.py}.



%% --------------------------------------------------------
\subsection{Висновки}
%% --------------------------------------------------------

\begin{itemize}
    \item Атом Гідрогену --- ідеальний тест для демонстрації роботи SCF-процедури в найпростішому випадку.
    \item Для одноелектронних систем HF є точним методом; похибка визначається лише базисом.
    \item Атом Гелію --- перша система, де з’являється електронна кореляція; HF дає наближення, але не точне значення енергії.
    \item Розширення базису покращує збіжність до повного базисного ліміту (CBS limit).
    \item Аналіз орбіталей і густини в PySCF дає інтуїцію про структуру розв’язку рівнянь Хартрі–Фока.
\end{itemize}

Таким чином, розгляд атомів H та He закладає основу для розуміння ітераційної природи SCF-процедури, ролі базисних функцій і меж застосовності методу Хартрі–Фока для багатoелектронних систем.
